\begin{abstract}
Nanosecond pulsed electric fields perturb cellular membranes and can alter extracellular vesicle release, but the intervening processes span incompatible timescales and are difficult to interpret from particle counts alone. Here we show a computational model that connects pulse delivery and membrane charging to calcium redistribution, bioenergetic stress, membrane repair, vesicle biogenesis and size-resolved extracellular particle dynamics. Simulations separated reversible, secretion-favored and injury-associated regimes. An observation model fitted particle distributions measured 0.5--3 h after exposure with a total-concentration error of 0.104 log$_{10}$ units and a mean distributional Hellinger distance of 0.086. The fit required an empirical exposure-dependent correction, indicating unresolved dose dependence, and apoptotic-body parameters remained weakly constrained. Sensitivity analysis prioritized pathway-to-particle conversion, basal release and extracellular loss. The model therefore provides a testable link between bioelectric exposure and heterogeneous particle measurements while identifying the measurements required for biological calibration.
\end{abstract}
